% ============================================================================
%  Data–MC Cell Reweighting Pipeline — Software Documentation
% ============================================================================
\documentclass[12pt,a4paper]{article}

% --- Packages ---------------------------------------------------------------
\usepackage[utf8]{inputenc}
\usepackage[T1]{fontenc}
\usepackage{lmodern}
\usepackage[margin=2.5cm,headheight=15pt]{geometry}
\usepackage{amsmath,amssymb}
\usepackage{graphicx}
\usepackage{booktabs}
\usepackage{caption}
\usepackage{hyperref}
\usepackage{cleveref}
\usepackage{xcolor}
\usepackage{float}
\usepackage{siunitx}
\usepackage{fancyhdr}
\usepackage{listings}
\usepackage{enumitem}
\usepackage{longtable}
\usepackage{array}

% --- Page style -------------------------------------------------------------
\pagestyle{fancy}
\fancyhf{}
\fancyhead[L]{\small Data--MC Reweighting Pipeline — Software Documentation}
\fancyhead[R]{\small ATLAS Internal}
\fancyfoot[C]{\thepage}
\renewcommand{\headrulewidth}{0.4pt}

% --- Hyperref setup ---------------------------------------------------------
\hypersetup{
  colorlinks=true,
  linkcolor=blue!60!black,
  citecolor=blue!60!black,
  urlcolor=blue!60!black
}

% --- Listings style ---------------------------------------------------------
\definecolor{codebg}{gray}{0.96}
\definecolor{codegreen}{rgb}{0.2,0.5,0.2}
\definecolor{codegray}{rgb}{0.5,0.5,0.5}

\lstdefinestyle{shellstyle}{
  backgroundcolor=\color{codebg},
  basicstyle=\ttfamily\small,
  breaklines=true,
  frame=single,
  rulecolor=\color{codegray},
  xleftmargin=4pt,
  xrightmargin=4pt,
  framesep=3pt,
  numbers=none,
  showstringspaces=false,
  commentstyle=\color{codegreen}\itshape,
  keywordstyle=\color{blue!70!black}\bfseries,
  stringstyle=\color{red!60!black},
  language=bash,
  morekeywords={root,bash,source,cd,mkdir},
  literate={~}{{\raisebox{0.5ex}{\texttildelow}}}{1},
}

\lstdefinestyle{cppstyle}{
  backgroundcolor=\color{codebg},
  basicstyle=\ttfamily\small,
  breaklines=true,
  frame=single,
  rulecolor=\color{codegray},
  xleftmargin=4pt,
  xrightmargin=4pt,
  framesep=3pt,
  numbers=left,
  numberstyle=\tiny\color{codegray},
  showstringspaces=false,
  commentstyle=\color{codegreen}\itshape,
  keywordstyle=\color{blue!70!black}\bfseries,
  stringstyle=\color{red!60!black},
  language=C++,
}

\lstset{style=shellstyle}

% --- Title ------------------------------------------------------------------
\title{%
  \textbf{Data--MC Cell Reweighting Pipeline}\\[6pt]
  \large Software Documentation\\[4pt]
  \large Photon Shower Shape Corrections — MC23e / Data24\\[4pt]
  \normalsize Version 2.0
}
\author{ATLAS Simulation}
\date{\today}

% ============================================================================
\begin{document}
\maketitle
\thispagestyle{fancy}

\tableofcontents
\clearpage

% ============================================================================
\section{Overview}
\label{sec:overview}
% ============================================================================

The data--MC cell reweighting pipeline derives and applies cell-level energy
corrections to photon shower shape variables in the ATLAS electromagnetic
calorimeter. The pipeline operates on ntuples produced by the
\texttt{NTupleMaker} Athena algorithm and produces corrected shower shape
distributions, validation plots, and diagnostic outputs.

The pipeline flow is:

\begin{enumerate}[leftmargin=2em]
  \item \textbf{Derive corrections} (\texttt{derive\_corrections.C}):
    Compute per-cell correction parameters (M1/M2/M3) from data and MC
    ntuples, stored in \texttt{corrections.root}.
  \item \textbf{Validate} (\texttt{validate\_data\_mc.C}):
    Apply corrections to MC, fill histograms of shower shape variables
    (branch-based and cell-computed), output to \texttt{histos.root}.
  \item \textbf{Plot shower shapes} (\texttt{plot\_data\_mc.C}):
    Overlay data vs MC (uncorrected and corrected) per $|\eta|$ bin.
  \item \textbf{Plot cell profiles} (\texttt{plot\_cell\_profiles.C}):
    Visualise $7\times11$ cell energy fraction maps and per-method residuals.
  \item \textbf{Plot stored vs computed} (\texttt{plot\_stored\_vs\_computed.C}):
    Compare branch-stored vs cell-recomputed shower shapes.
\end{enumerate}

All scripts share a common configuration header (\texttt{config.h}) and
support parameterised isolation levels for systematic studies.

% ============================================================================
\section{Software Requirements}
\label{sec:requirements}
% ============================================================================

\begin{table}[htbp]
\centering
\caption{Software dependencies.}
\label{tab:requirements}
\begin{tabular}{llp{8cm}}
\toprule
\textbf{Software} & \textbf{Version} & \textbf{Purpose} \\
\midrule
ROOT & 6.28.04 & Event processing, histogramming, TProfile, plotting \\
OS   & AlmaLinux~9 (EL9) & CERN lxplus cluster \\
CVMFS & \texttt{/cvmfs/atlas.cern.ch/} & Software distribution \\
\bottomrule
\end{tabular}
\end{table}

\subsection{Environment setup}

\begin{lstlisting}
export ATLAS_LOCAL_ROOT_BASE=/cvmfs/atlas.cern.ch/repo/ATLASLocalRootBase
source ${ATLAS_LOCAL_ROOT_BASE}/user/atlasLocalSetup.sh --quiet
lsetup "root 6.28.04-x86_64-el9-gcc13-opt" --quiet
\end{lstlisting}

% ============================================================================
\section{Repository Structure}
\label{sec:structure}
% ============================================================================

\begin{longtable}{p{5.5cm}p{9cm}}
\toprule
\textbf{File} & \textbf{Description} \\
\midrule
\endhead
\texttt{scripts/data\_mc/config.h} & Shared configuration: branch names, $|\eta|$ binning, cluster geometry, selection cuts, isolation parameters, helper functions \\
\texttt{scripts/data\_mc/derive\_corrections.C} & Step~1: Derive M1/M2/M3 corrections from data+MC \\
\texttt{scripts/data\_mc/validate\_data\_mc.C} & Step~2: Apply corrections, fill histograms \\
\texttt{scripts/data\_mc/plot\_data\_mc.C} & Step~3: Overlay shower shape distributions \\
\texttt{scripts/data\_mc/plot\_cell\_profiles.C} & Step~4: Cell profile maps + per-method residuals \\
\texttt{scripts/data\_mc/plot\_stored\_vs\_computed.C} & Step~5: Branch vs cell-computed comparison \\
\texttt{scripts/data\_mc/plot\_fudge\_binned.C} & Fudge factor validation (TUNE27AD binning) \\
\texttt{scripts/data\_mc/run\_pipeline.sh} & Run full pipeline for one channel + isolation \\
\texttt{scripts/data\_mc/run\_all\_channels.sh} & Run all channels $\times$ isolation variants \\
\bottomrule
\end{longtable}

% ============================================================================
\section{Configuration (\texttt{config.h})}
\label{sec:config}
% ============================================================================

All analysis parameters are centralised in \texttt{config.h} within the
\texttt{datamc} namespace.

\subsection{Cluster geometry}

\begin{lstlisting}[style=cppstyle]
const int kPhiSize = 11;  // phi direction
const int kEtaSize = 7;   // eta direction
const int kClusterSize = kPhiSize * kEtaSize;  // 77 cells
\end{lstlisting}

\subsection{Pseudorapidity binning}

14~bins following the standard ATLAS EM calorimeter segmentation.
Bin~8 ($1.37 \leq |\eta| < 1.52$) is the barrel--endcap crack region
and is always empty due to passEtaAcceptance rejection.

\subsection{Energy fraction binning (M2)}

9~bins in the physical range $[0, 0.5]$ for M2 TProfile histograms:
\begin{lstlisting}[style=cppstyle]
const double kFracBins[] = {0.0, 0.02, 0.04, 0.06, 0.08,
                            0.10, 0.15, 0.20, 0.30, 0.50};
const int kNFracBins = 9;
\end{lstlisting}

\subsection{Isolation levels}

\begin{lstlisting}[style=cppstyle]
// Thresholds
const float kPtcone20OverPtMax = 0.05;
const float kTopoetcone40OverPtMax = 2.472;

// Sigma ratio capping for M3
const double kSigmaRatioMin = 0.5;
const double kSigmaRatioMax = 2.0;
\end{lstlisting}

The \texttt{passSelection()} function accepts an \texttt{isoLevel}
parameter (0--3) that gates the isolation requirements:
\begin{itemize}
  \item \texttt{isoLevel $\geq$ 1}: Require \texttt{isotight} $= 1$.
  \item \texttt{isoLevel $\geq$ 2}: Additionally require
    \texttt{ptcone20}$/\pt < 0.05$.
  \item \texttt{isoLevel $\geq$ 3}: Additionally require
    \texttt{topoetcone40}$/\pt < 2.472$.
\end{itemize}

All isolation parameters have default values of zero, maintaining
backward compatibility with existing call sites.

% ============================================================================
\section{Script Reference}
\label{sec:scripts}
% ============================================================================

\subsection{\texttt{derive\_corrections.C}}
\label{sec:scripts:derive}

\paragraph{Purpose.} Derive per-cell correction parameters for all three
methods from data and MC ntuples.

\paragraph{Signature.}
\begin{lstlisting}[style=cppstyle]
void derive_corrections(const char* dataFile,
                        const char* mcFile,
                        const char* outFile,
                        Long64_t maxEvents = -1,
                        int isoLevel = 0);
\end{lstlisting}

\paragraph{Output.} \texttt{corrections.root} containing:
\begin{itemize}
  \item \texttt{cellsMeanData\_Eta\_LO\_HI} (TH2D): Mean data fractions.
  \item \texttt{cellsMeanMC\_Eta\_LO\_HI} (TH2D): Mean MC fractions.
  \item \texttt{cellsDelta\_Eta\_LO\_HI} (TH2D): M1 shifts.
  \item \texttt{cellsSigmaRatio\_Eta\_LO\_HI} (TH2D): M3 sigma ratios
    (capped to $[0.5, 2.0]$).
  \item \texttt{m2\_alpha\_Cell\_K\_Eta\_LO\_HI} (TProfile): M2 profiled
    corrections per cell per $|\eta|$ bin.
\end{itemize}

\subsection{\texttt{validate\_data\_mc.C}}
\label{sec:scripts:validate}

\paragraph{Purpose.} Apply corrections to MC events and fill shower shape
histograms for data and MC (corrected and uncorrected).

\paragraph{Signature.}
\begin{lstlisting}[style=cppstyle]
int validate_data_mc(const char* dataFile,
                     const char* mcFile,
                     const char* corrFile,
                     const char* histoFile,
                     Long64_t maxEvents = -1,
                     int isoLevel = 0);
\end{lstlisting}

\paragraph{Output.} \texttt{histos.root} containing per-$|\eta|$-bin
histograms:
\begin{itemize}
  \item Set~A (branch-based): \texttt{\{reta,rphi,weta2\}\_\{data\_br,mc\_unf,mc\_fud\}\_N}
  \item Set~B (cell-computed): \texttt{\{reta,rphi,weta2\}\_\{data\_cell,mc\_cell,m1,m2,m3\}\_N}
\end{itemize}

\subsection{\texttt{plot\_data\_mc.C}}
\label{sec:scripts:plot_data_mc}

\paragraph{Purpose.} Produce overlay plots of data vs MC (corrected and
uncorrected) per $|\eta|$ bin.

\paragraph{Signature.}
\begin{lstlisting}[style=cppstyle]
int plot_data_mc(const char* histoFile, const char* plotDir);
\end{lstlisting}

\paragraph{Output.} Multi-page PDFs:
\texttt{overlay\_\{reta,rphi,weta2\}.pdf}.

\subsection{\texttt{plot\_cell\_profiles.C}}
\label{sec:scripts:cell_profiles}

\paragraph{Purpose.} Visualise $7\times11$ cell energy fraction maps and
per-method correction residuals.

\paragraph{Signature.}
\begin{lstlisting}[style=cppstyle]
int plot_cell_profiles(const char* corrFile,
                       const char* plotDir,
                       const char* channelLabel = "...");
\end{lstlisting}

\paragraph{Output.}
\begin{itemize}
  \item \texttt{cell\_profiles\_before.pdf}: Data vs uncorrected MC.
  \item \texttt{cell\_profiles\_after\_m1.pdf}: Data vs M1-corrected MC.
  \item \texttt{cell\_profiles\_after\_m2.pdf}: Data vs M2-corrected MC.
  \item \texttt{cell\_profiles\_after\_m3.pdf}: Data vs M3-corrected MC.
  \item \texttt{cell\_profiles\_delta.pdf}: Delta maps and sigma ratios.
  \item \texttt{cell\_delta\_m1.pdf}: M1 residual
    ($\langle f\rangle_\text{data} - \langle f\rangle_\text{M1}$).
  \item \texttt{cell\_delta\_m2.pdf}: M2 residual.
  \item \texttt{cell\_delta\_m3.pdf}: M3 residual.
\end{itemize}

\subsection{\texttt{plot\_stored\_vs\_computed.C}}
\label{sec:scripts:stored_vs_computed}

\paragraph{Purpose.} Overlay branch-stored shower shapes against
cell-recomputed values to validate the computation.

\paragraph{Signature.}
\begin{lstlisting}[style=cppstyle]
int plot_stored_vs_computed(const char* histoFile,
                            const char* plotDir,
                            const char* channelLabel = "...");
\end{lstlisting}

\paragraph{Output.} Multi-page PDFs:
\texttt{stored\_vs\_computed\_\{reta,rphi,weta2\}.pdf}.

\subsection{\texttt{plot\_fudge\_binned.C}}
\label{sec:scripts:fudge}

\paragraph{Purpose.} Validate standard fudge factors using TUNE27AD
$\eta\times p_\text{T}$ binning.

\paragraph{Signature.}
\begin{lstlisting}[style=cppstyle]
void plot_fudge_binned(Long64_t maxEv = -1,
                       int isoLevel = 0);
\end{lstlisting}

\paragraph{Output.} PDFs in
\texttt{data\_mc\_reweighting\_\{ISO\_TAG\}/\{channel\}/}:
\texttt{fudge\_inclusive.pdf}, \texttt{fudge\_etabins.pdf},
\texttt{fudge\_eta\_pt.pdf}, plus
\texttt{fudge\_comparison\_\{reta,rphi,weta2\}.pdf}.

% ============================================================================
\section{Output Directory Structure}
\label{sec:output}
% ============================================================================

Output is organised by isolation tag and channel:
\begin{lstlisting}
output/
  data_mc_reweighting_{ISO_TAG}/
    Zeeg/
      corrections.root
      histos.root
      plots/
        overlay_reta.pdf
        overlay_rphi.pdf
        overlay_weta2.pdf
        cell_profiles_before.pdf
        cell_profiles_after_m1.pdf
        cell_profiles_after_m2.pdf
        cell_profiles_after_m3.pdf
        cell_profiles_delta.pdf
        cell_delta_m1.pdf
        cell_delta_m2.pdf
        cell_delta_m3.pdf
        stored_vs_computed_reta.pdf
        stored_vs_computed_rphi.pdf
        stored_vs_computed_weta2.pdf
    Zmumug/
      ...
    combined/
      ...
    fudge_inclusive.pdf
    fudge_etabins.pdf
    ...
\end{lstlisting}

Where \texttt{ISO\_TAG} is one of: \texttt{no\_iso}, \texttt{iso\_tight},
\texttt{iso\_track}, \texttt{iso\_full}.

% ============================================================================
\section{Running the Analysis}
\label{sec:running}
% ============================================================================

\subsection{Single channel + isolation variant}

\begin{lstlisting}
cd scripts/data_mc/
bash run_pipeline.sh Zeeg -1 iso_full 3
\end{lstlisting}

Arguments:
\begin{enumerate}
  \item Channel: \texttt{Zeeg}, \texttt{Zmumug}, or \texttt{combined}.
  \item Max events: \texttt{-1} for all events.
  \item Isolation tag: \texttt{no\_iso}, \texttt{iso\_tight},
    \texttt{iso\_track}, \texttt{iso\_full}.
  \item Isolation level: \texttt{0}--\texttt{3}, must match the tag.
\end{enumerate}

\subsection{All channels and isolation variants}

\begin{lstlisting}
bash run_all_channels.sh all -1
\end{lstlisting}

This runs all 12~combinations (4~isolation $\times$ 3~channels) plus the
fudge factor validation for each isolation level.

\subsection{Single isolation variant only}

\begin{lstlisting}
bash run_all_channels.sh iso_full -1
\end{lstlisting}

\subsection{Quick test with limited events}

\begin{lstlisting}
bash run_all_channels.sh no_iso 50000
\end{lstlisting}

% ============================================================================
\section{Reproducibility}
\label{sec:reproducibility}
% ============================================================================

The analysis is fully deterministic given the same input ntuples. To
reproduce all results:

\begin{enumerate}
  \item Ensure access to the input ntuples at
    \texttt{/dcache/atlas/mfernand/qt\_ntuples/data24/}.
  \item Set up the ROOT environment as described in \cref{sec:requirements}.
  \item Navigate to \texttt{scripts/data\_mc/} and run
    \texttt{bash run\_all\_channels.sh all -1}.
  \item Output appears in \texttt{output/data\_mc\_reweighting\_\{ISO\_TAG\}/}.
\end{enumerate}

All configuration parameters (cuts, binning, isolation thresholds, sigma
ratio caps) are defined in \texttt{config.h}. Changing any parameter and
rerunning the pipeline produces a new, self-consistent set of results.

\end{document}
